\documentclass[10pt,conference,compsocconf]{IEEEtran}

\usepackage{amsmath}
\usepackage{amssymb}
\usepackage{graphicx}
\usepackage{hyperref}

\begin{document}
\title{Application of Machine Learning in Predictive Coronary Heart Disease (MICHD) Diagnosis}

\author{Paul Barbereau, Amandine Dejardin, Samuel Perry\\
\textit{School of Computer and Communication Sciences, EPFL, Lausanne, Switzerland}}

\maketitle

\begin{abstract}
Cardiovascular diseases (CVDs) are a leading cause of mortality worldwide, making early risk screening from clinical and lifestyle indicators a key public health goal. We present a machine learning pipeline to predict myocardial infarction or coronary heart disease (MICHD) from the 2015 Behavioral Risk Factor Surveillance System (BRFSS) dataset ($328{,}135$ training and $109{,}379$ test individuals across $321$ variables). To address missing data, survey-specific coding, feature redundancy, and severe class imbalance ($8.83\%$ positive rate), we combine domain-informed selection of $47$ clinically relevant features with codebook-guided recoding, fold-wise median imputation, standardization, and $\ell_2$-regularized logistic regression with threshold tuning. In 5-fold cross-validation, our model achieves an F1 score of $0.4199 \pm 0.0030$ and $87.15\% \pm 0.05\%$ accuracy, substantially outperforming uncleaned baselines and linear regression models.
\end{abstract}

\section{Introduction}
Early identification of individuals at risk of myocardial infarction or coronary heart disease (MICHD) using non-invasive health surveys can enable timely preventive care. In this work, we predict binary MICHD diagnosis ($y \in \{-1, 1\}$) from the CDC's 2015 BRFSS telephone survey~\cite{cdc_brfss_codebook_2015}. Applying machine learning directly to the raw dataset presents four main challenges: (1) \textbf{high dimensionality} ($D = 321$ columns, including administrative metadata unrelated to health); (2) \textbf{survey-specific sentinel codes}, where ``None/Zero,'' ``Don't know,'' and ``Refused'' are encoded as large integers ($7, 9, 77, 88, 99, 888, 9999$) that distort numerical distributions; (3) \textbf{multicollinearity} among raw and derived variables; and (4) \textbf{class imbalance}, as only $28{,}975$ of $328{,}135$ training respondents ($8.83\%$) report MICHD. To address these issues without external libraries, we implement six foundational algorithms in \texttt{implementations.py} and pair $\ell_2$-regularized logistic regression with codebook-driven preprocessing and 5-fold cross-validation in \texttt{run.py}.

\section{Models and Methodology}

\subsection{Model Selection and Baselines}
We implemented the six core methods in \texttt{implementations.py}: linear models minimizing mean squared error (\texttt{least\_squares}, \texttt{ridge\_regression}, \texttt{mean\_squared\_error\_gd}, and mini-batch $B=1$ \texttt{mean\_squared\_error\_sgd}) and probabilistic classifiers minimizing negative log-likelihood (\texttt{logistic\_regression} and \texttt{reg\_logistic\_regression}). Because MICHD diagnosis is binary, we selected $\ell_2$-regularized logistic regression (with $\lambda = 0$ corresponding to standard logistic regression). Given a feature vector $\mathbf{x}_i \in \mathbb{R}^{D^*}$ and sigmoid $\sigma(z) = (1 + e^{-z})^{-1}$, the positive class probability is $p(y_i = 1 \mid \mathbf{x}_i) = \sigma(\mathbf{w}^T \mathbf{x}_i)$, and the regularized training loss is:
\begin{equation}
\begin{aligned}
    \mathcal{L}(\mathbf{w}) = -\frac{1}{N} \sum_{i=1}^{N} \Big[ &y_i \log\sigma(\mathbf{w}^T \mathbf{x}_i) \\
    &+ (1 - y_i) \log\big(1 - \sigma(\mathbf{w}^T \mathbf{x}_i)\big) \Big] + \lambda \|\mathbf{w}\|_2^2,
\end{aligned}
\end{equation}
where $\lambda \geq 0$ penalizes large weights to prevent overfitting.

\subsection{Data Preprocessing}
Our preprocessing pipeline in \texttt{run.py} transforms the $321$ raw columns into a clean design matrix of $D^* = 48$ features ($47$ variables plus an intercept):
\begin{enumerate}
    \item \textbf{Feature Selection:} Using the BRFSS codebook~\cite{noauthor_www2statdukeeducr173sta102_sp16projcodebookhtml_nodate,cdc_brfss_codebook_2015}, we selected $47$ variables (\texttt{cvd\_variables}) covering general health (\texttt{GENHLTH}, \texttt{PHYSHLTH}), cardiovascular history (\texttt{BPHIGH4}, \texttt{TOLDHI2}, \texttt{CVDSTRK3}, \texttt{HAREHAB1}), diabetes (\texttt{DIABETE3}), demographics (\texttt{SEX}, \texttt{\_AGE80}, \texttt{EDUCA}, \texttt{INCOME2}), and lifestyle factors (smoking, alcohol, exercise, BMI).
    \item \textbf{Survey Code Recoding:} In \texttt{preprocess\_cols}, codes representing ``None'' or zero days/drinks ($88, 888$) are mapped to $0$, while ``Don't know,'' ``Refused,'' or missing codes ($7, 9, 77, 98, 99, 777, 999, 7777, 9999, 3, 14$) are set to \texttt{NaN}.
    \item \textbf{Imputation and Standardization:} In \texttt{preprocessing\_X}, all-\texttt{NaN} and constant columns are removed, and remaining \texttt{NaN} values are imputed with the training column median. Features are standardized using the training mean and standard deviation to ensure stable gradient descent, a bias column of ones is added, and labels $y \in \{-1, 1\}$ are mapped to $\{0, 1\}$.
\end{enumerate}

\subsection{Optimization, Hyperparameter Tuning, and Prediction}
Weights $\mathbf{w}$ are optimized over $1{,}000$ gradient descent iterations via $\mathbf{w} \leftarrow \mathbf{w} - \gamma \nabla \mathcal{L}(\mathbf{w})$, where:
\begin{equation}
    \nabla \mathcal{L}(\mathbf{w}) = \frac{1}{N} \sum_{i=1}^{N} \big(\sigma(\mathbf{w}^T \mathbf{x}_i) - y_i\big) \mathbf{x}_i + 2\lambda \mathbf{w}.
\end{equation}
Test predictions are obtained by thresholding the predicted probability at $\tau$: $\hat{y}_i = 1$ if $\sigma(\mathbf{w}^T \mathbf{x}_i) \geq \tau$ and $-1$ (or $0$) otherwise. We tuned $(\lambda, \gamma, \tau)$ via 5-fold cross-validation (\texttt{SEED = 67}) maximizing the F1 score:
\begin{equation}
    \text{F1} = \frac{2 \cdot \text{Precision} \cdot \text{Recall}}{\text{Precision} + \text{Recall}},
\end{equation}
with $\text{Precision} = \frac{\text{TP}}{\text{TP}+\text{FP}}$ and $\text{Recall} = \frac{\text{TP}}{\text{TP}+\text{FN}}$. Imputation and standardization were computed inside each fold to prevent data leakage. Grid search over $\lambda \in \{0\} \cup [10^{-6}, 10^0]$, $\gamma \in [10^{-3}, 10^0]$, and $\tau \in [0.1, 0.7]$ yielded $\lambda = 10^{-6}$, $\gamma = 0.775$, and $\tau = 0.20$.

\section{Results}

\subsection{Comparison with Baseline Models}
Table~\ref{tab:baselines} compares all six implemented algorithms on the cleaned 47-feature dataset across 5 cross-validation folds. Unregularized \texttt{least\_squares} fails with a singular matrix because derived BRFSS variables (e.g., \texttt{\_AGEG5YR}, \texttt{\_AGE65YR}, \texttt{\_AGE80}, \texttt{\_AGE\_G}) are collinear, making $\mathbf{X}^T\mathbf{X}$ rank-deficient. Adding an $\ell_2$ penalty in \texttt{ridge\_regression} ($\lambda = 10^{-6}$) restores invertibility and matches full-batch \texttt{mse\_gd} ($\text{F1} = 0.4109$ at $\tau = 0.20$). Meanwhile, \texttt{mse\_sgd} ($B=1$, $1{,}000$ steps) observes only $0.38\%$ of each fold and suffers from high variance ($\text{F1} = 0.3106 \pm 0.0852$). Both \texttt{logistic\_regression} and \texttt{reg\_logistic\_regression} at $\tau = 0.20$ outperform all linear MSE baselines ($\text{F1} = 0.4199$ vs.\ $0.4109$, accuracy $87.15\%$ vs.\ $85.76\%$), confirming that the log-loss objective and bounded sigmoid output better capture binary disease risk.

\begin{table}[htbp]
\centering
\caption{5-Fold CV Comparison on Cleaned 47 Features ($\text{Mean} \pm \text{SD}$)}
\label{tab:baselines}
\resizebox{\columnwidth}{!}{%
\begin{tabular}{|l|c|c|c|c|c|}
\hline
\textbf{Method} & $\boldsymbol{\tau}$ & \textbf{F1 Score} & \textbf{Accuracy} & \textbf{Prec.} & \textbf{Rec.} \\
\hline
\texttt{least\_squares} & -- & \multicolumn{4}{c|}{\textit{Singular matrix ($\sigma_{\min} = 2.6 \times 10^{-10}$)}} \\
\hline
\texttt{ridge\_reg} ($\lambda\!=\!10^{-6}$) & $0.50$ & $0.0766 \pm 0.0021$ & $0.9140 \pm 0.0011$ & $0.734$ & $0.040$ \\
\texttt{ridge\_reg} ($\lambda\!=\!10^{-6}$) & $0.20$ & $0.4109 \pm 0.0032$ & $0.8576 \pm 0.0010$ & $0.324$ & $0.563$ \\
\hline
\texttt{mse\_gd} ($\gamma\!=\!0.1$) & $0.50$ & $0.0767 \pm 0.0021$ & $0.9140 \pm 0.0011$ & $0.733$ & $0.041$ \\
\texttt{mse\_gd} ($\gamma\!=\!0.1$) & $0.20$ & $0.4109 \pm 0.0030$ & $0.8576 \pm 0.0010$ & $0.324$ & $0.562$ \\
\hline
\texttt{mse\_sgd} ($\gamma\!=\!10^{-3}$) & $0.20$ & $0.3106 \pm 0.0852$ & $0.8836 \pm 0.0196$ & $0.367$ & $0.328$ \\
\hline
\texttt{logistic\_reg} & $0.50$ & $0.2185 \pm 0.0098$ & $\mathbf{0.9146 \pm 0.0012}$ & $0.568$ & $0.135$ \\
\texttt{logistic\_reg} & $0.20$ & $\mathbf{0.4199 \pm 0.0030}$ & $0.8715 \pm 0.0005$ & $0.349$ & $0.527$ \\
\hline
\textbf{\texttt{reg\_log\_reg}} ($\lambda\!=\!10^{-6}$) & $0.50$ & $0.2185 \pm 0.0098$ & $\mathbf{0.9146 \pm 0.0012}$ & $0.568$ & $0.135$ \\
\textbf{\texttt{reg\_log\_reg}} ($\lambda\!=\!10^{-6}$) & $\mathbf{0.20}$ & $\mathbf{0.4199 \pm 0.0030}$ & $0.8715 \pm 0.0005$ & $0.349$ & $0.527$ \\
\hline
\end{tabular}%
}
\end{table}

\subsection{Ablation Study}
Table~\ref{tab:ablation} isolates the impact of each pipeline modification. Compared to using all $315$ non-constant features (Config.\ 1), restricting to our $47$ domain features (Config.\ 2) cuts dimensionality by $85\%$ and training time by $6.5\times$. Adding BRFSS survey code cleaning (Config.\ 3) lifts F1 at $\tau = 0.50$ from $0.1765$ to $0.2185$ ($+4.20$ points). Finally, tuning the threshold from $\tau = 0.50$ to $\tau = 0.20$ (Config.\ 6) provides a $+20.14$ point F1 jump ($0.2185 \to 0.4199$), outperforming the 315-feature model at $\tau = 0.20$ (Config.\ 5, $0.4149$) in both F1 and accuracy.

\begin{table}[htbp]
\centering
\caption{Ablation Study of Pipeline Modifications (5-Fold CV)}
\label{tab:ablation}
\resizebox{\columnwidth}{!}{%
\begin{tabular}{|l|c|c|c|}
\hline
\textbf{Pipeline Configuration} & $\boldsymbol{D^*}$ & \textbf{CV F1 Score} & \textbf{CV Accuracy} \\
\hline
1. All raw features, no code clean, $\tau=0.50$ & $316$ & $0.1814 \pm 0.0067$ & $0.9142 \pm 0.0007$ \\
2. 47 features, no code clean, $\tau=0.50$ & $48$ & $0.1765 \pm 0.0068$ & $0.9135 \pm 0.0008$ \\
3. 47 features + code clean, $\tau=0.50$ & $48$ & $0.2185 \pm 0.0098$ & $\mathbf{0.9146 \pm 0.0012}$ \\
4. 47 features, no code clean, $\tau=0.20$ & $48$ & $0.4113 \pm 0.0039$ & $0.8671 \pm 0.0009$ \\
5. All raw features, no code clean, $\tau=0.20$ & $316$ & $0.4149 \pm 0.0020$ & $0.8659 \pm 0.0005$ \\
\hline
\textbf{6. Full Pipeline (47 + clean, $\boldsymbol{\tau=0.20}$)} & $\mathbf{48}$ & $\mathbf{0.4199 \pm 0.0030}$ & $0.8715 \pm 0.0005$ \\
\hline
\end{tabular}%
}
\end{table}

\section{Discussion}
\textbf{Class Imbalance and Thresholding:} Because only $8.83\%$ of patients have MICHD, predicting $-1$ everywhere yields $91.17\%$ accuracy but $\text{F1} = 0$. At $\tau = 0.50$, logistic regression achieves $91.46\%$ accuracy but only $13.53\%$ recall ($\text{F1} = 0.2185$). Lowering $\tau$ to $0.20$ accounts for the low positive prior, nearly quadrupling recall to $52.69\%$ and doubling F1 to $0.4199$. Clinically, higher recall is critical in screening, where missed heart disease (false negatives) is far costlier than follow-up testing (false positives).

\textbf{Clinical Interpretability and Limitations:} On the standardized training set (final log-loss $0.2231$, intercept $w_0 = -3.254$), the strongest predictors of MICHD are poor general health (\texttt{GENHLTH}, $w = +0.455$), older age (\texttt{\_AGE80}, $w = +0.442$; \texttt{\_AGEG5YR}, $w = +0.405$), hypertension (\texttt{\_RFHYPE5}, $w = +0.142$), high cholesterol (\texttt{\_RFCHOL}, $w = +0.137$), male sex (\texttt{SEX}, $w = -0.359$), prior heart rehabilitation (\texttt{HAREHAB1}, $w = -0.262$), and stroke (\texttt{CVDSTRK3}, $w = -0.184$), matching established medical risk factors. Limitations include standardizing categorical variables without one-hot encoding, retaining collinear age features, and assuming a linear boundary. Future work could add one-hot encoding, $\ell_1$ feature selection, polynomial interactions, and class-weighted loss.

\begin{thebibliography}{99}
\bibitem{noauthor_www2statdukeeducr173sta102_sp16projcodebookhtml_nodate}
Duke University, ``BRFSS 2015 Project Codebook,'' STA 102 Course Materials, \url{https://www2.stat.duke.edu/~cr173/Sta102_Sp16/Proj/codebook.html}.
\bibitem{cdc_brfss_codebook_2015}
CDC, ``Behavioral Risk Factor Surveillance System 2015 Codebook Report,'' 2016.
\end{thebibliography}

\end{document}
